% =============================================================================
% SECTION 1: APPLICATION AREA
% =============================================================================

\section{Application Area}

\subsection{Description}

The application for this semester project is \textbf{structural feasibility assessment of wood-framed residential headers and beams} within a construction automation system called Construction.AI.

In residential wood-frame construction, horizontal members---headers spanning door and window openings, ridge beams, and floor beams---must be sized to carry tributary loads from the roof and floors above without exceeding allowable stress or deflection limits prescribed by the International Residential Code (IRC).
Structural engineers and experienced estimators perform this sizing manually using prescriptive span tables or first-principles calculations.
The Construction.AI project seeks to automate this process: given a set of architectural floor plans, the system reads the geometry, identifies load-bearing walls and openings, generates plausible structural hypotheses, and evaluates each hypothesis by solving the governing structural mechanics equations numerically.

The core physics model governing each header or beam is the \textbf{Euler--Bernoulli beam equation}, a fourth-order ordinary differential equation (ODE) relating beam deflection to the applied distributed load.
This is a 1D steady boundary value problem (BVP) whose independent variable is position $x$ along the beam axis---a simple, analytically tractable PDE that satisfies the course project criteria.

\subsection{Why This Application}

This application satisfies all recommended project criteria:

\begin{enumerate}
  \item \textbf{Simple physics.} The Euler--Bernoulli beam model is a 1D linear fourth-order ODE, one of the most classical PDEs in structural mechanics.
  \item \textbf{One independent variable.} The spatial coordinate $x \in [0, L]$ is the sole independent variable; the problem is steady-state.
  \item \textbf{Analytical solution available.} Closed-form solutions exist for canonical loading cases (uniform load, point load, simply-supported and fixed-fixed boundary conditions), enabling rigorous code verification without manufactured solutions.
  \item \textbf{No singularities.} Under smooth distributed loads the solution $w(x)$ is $C^4$ smooth, with no internal singularities.
  \item \textbf{Software already written.} A finite-difference beam solver has been designed as part of the Construction.AI system; the source code is accessible and modifiable.
  \item \textbf{Research relevance.} The project is directly tied to ongoing work automating construction plan analysis---a practical, industry-relevant application.
  \item \textbf{Tractable for uncertainty propagation.} Each beam solve completes in milliseconds, making Monte Carlo runs of $10^3$--$10^4$ samples computationally cheap.
\end{enumerate}

\subsection{Team}

This assignment is completed as a \textbf{2-person team}:
\begin{itemize}
  \item Jason Cusati (Virginia Tech)
  \item Cheng-Shun Chuang (Virginia Tech)
\end{itemize}